% !TeX root = ../thuthesis-example.tex

\chapter{国内外相关工作}

\section{梯度压缩算法}

\subsection{非自适应的传统梯度压缩算法}

根据梯度压缩综述 \cite{grace,孙志军2012深度学习研究综述}，传统梯度压缩算法可以被分为稀疏化、量化、低秩压缩算法以及他们的混合（如图\ref{fig:TraditionCompression}）。稀疏化压缩算法只传递梯度中的部分值，其他的都设置为零 \cite{alistarh2018convergence,hardthreshold,stich2019error}，Han Song等人在梯度稀疏化的基础上采用动量修正、局部梯度裁剪等技巧保证模型质量，最高可以实现1/600的压缩率\cite{dgc}；量化压缩算法在保证梯度数学期望不变的前提下，将高精度（如32位浮点数）映射成低精度（如4位整数） \cite {bernstein2018signsgd,wu2018error,qsgd}；低秩算法是将一个梯度矩阵解耦成两个或者多个低秩矩阵。Thijs Vogels等人通过低秩压缩设计出了PowerSGD，该算法可以快速压缩梯度，在各个优化后端上都实现了和SGD一致的时钟加速\cite{vogels2019powersgd}。混合梯度压缩算法是指
将多种梯度压缩策略混合的梯度压缩算法\cite{qsparse,stc,cui2021slashing}，有工作提出了稀疏化二元梯度压缩算法（如图\ref{fig:stc}）。他们在原先SBC算法的基础上，采用了下游压缩和权重更新缓存等策略增强了算法的鲁棒性。虽然传统的分布式机器学习中，梯度压缩已经取得了卓越的成就，但是严重Non-IID问题会让梯度压缩算法出现精度损失  \cite{niid2020icml}，设计对Non-IID高度鲁棒性的梯度压缩算法对梯度压缩算法在新型分布式机器学习中的应用有十分重要的意义。

\begin{figure}
  \centering
  \includegraphics[width=1.0\linewidth]{figures/2-1.pdf}
  %\caption*{1954年至2021年间流行的新机器学习系统的模型规模。}
  \caption{传统梯度压缩算法可以被分为量化（Sign，Sign+Norm）、稀疏化（Top-K，Random-K）和低秩梯度压缩算法（Low-rank）\cite{vogels2019powersgd}。}
  \label{fig:TraditionCompression}
\end{figure}

\begin{figure}
  \centering
  \includegraphics[width=1.0\linewidth]{figures/stc.pdf}
  %\caption*{1954年至2021年间流行的新机器学习系统的模型规模。}
  \caption{SBC算法示意图 \cite{sbc}。}
  \label{fig:stc}
\end{figure}


相对压缩器在稀疏化压缩器中非常流行\cite{stich2020communication,stich2019error,aji2017sparse}。 我们把压缩率表示为$\delta$，把相对压缩器表示为$\textbf{C}_{\delta}$。$\textbf{C}_{\delta}$是一个映射：$\mathbb{R}^d \rightarrow \mathbb{R}^d$，其数学特征是：
\begin{equation}
    \mathbb{E}_{\textbf{C}_{\delta}}\lVert \textbf{C}_{\delta}(\textbf{x}) - \textbf{x} \rVert^2 \leq (1-\delta) \lVert \textbf{x} \rVert^2 \nonumber
\end{equation}

最近有关于绝对压缩器的工作 \cite{hardthreshold}，它对误差有一个绝对的约束，但在我们并不考虑这种情况。
：采用任意无偏量化器的分布式随机梯度下降（以下简称Q-SGD）和采用带误差反馈机制的有偏压缩器的随机梯度下降（以下简称D-EF-SGD）\cite{cui2021slashing,QuantityGlobal}是新型分布式机器学习中的两种主流的压缩器。前者在通信受限的情况下有两个优势。第一个优势是，具有相对压缩器的D-EF-SGD可以实现较低的压缩比\footnote{D-QSGD的最小压缩比可以达到$\frac{1}{32}$\cite{signsgd,qsgd}，接近 $3.4\%$，而稀疏化梯度压缩的压缩比可以低于$0.1\%$ \cite{aji2017sparse}。更重要的是，D-QSGD的压缩比是离散的，这使得自适应调节问题更加复杂，灵活度更低。这个问题在DC2 \cite{dc2}的设计中也有讨论。} （即传输更少的数据），它更适合于资源受限的情况。第二个优点是D-EF-SGD对数据异质性$\zeta$的依赖性较小，更适合于高偏度的Non-IID分布情况\cite{stich2020communication}。因此，在这项工作中主要考虑了具有相对压缩器的D-EF-SGD。


 
\subsection{自适应梯度压缩算法}

自适应梯度压缩算法可以自适应的调整传统压缩算法中的压缩参数。这有助于增强梯度压缩算法对一些特定场景的健壮性，如动态网络场景和有损压缩自带的噪音 \cite{AdaQS,PASGD}。有工作\cite{accordion}提出了Accordion算法来识别当前的训练阶段是否是关键阶段。如果是关键阶段就采用保守压缩来传输更多的数据，否则就采用激进压缩并降低传输数据量，提升训练速度，这从时间维度更细粒度的考虑了压缩率的变化，一定程度上增加了传统梯度压缩算法的鲁棒性，但是该算法更多是一种经验上的考量，缺乏理论论证过程，不能说明自适应方法的普遍适用性。也有工作\cite{dc2}提出了名为DC2的基于网络延迟的压缩控制系统（见图\ref{fig:2.2}）。DC2通过监测网络延迟变化来更新压缩参数，在网络传输状况变差或者可能变差时，通过采用激进的压缩率设置来有效应对网络问题，在网络状况恢复常态时采用保守的压缩率设置，从而有效地增强了梯度压缩算法对动态网络环境的鲁棒性，加速了动态网络环境下的模型收敛，但是该算法有两个不足，
一是引入了较多的超参数，另一方面缺乏理论证明，缺少算法的收敛性保证。还有一些工作从提升梯度压缩算法对非IID场景的鲁棒性角度，提出了一些自适应梯度压缩算法。有工作\cite{niid2020icml}提出了SkewScout策略 （见图\ref{fig:2.3}）。这种系统级的方法根据节点之间的损失（或精确值）的差异来动态地调整梯度压缩超参数。然而，在实践中，我们很难用一个同一只的指标量化各节点数据分布的差异，这使得SkewScout难以落地。

\begin{figure}
  \centering
  \includegraphics[width=0.8\linewidth]{figures/2-2.pdf}
  \caption{DC2的示意图\cite{dc2}。}
  \label{fig:2.2}
\end{figure}

\begin{figure}
  \centering
  \includegraphics[width=0.8\linewidth]{figures/2-3.pdf}
  \caption{SkewScout的示意图\cite{dc2}。}
  \label{fig:2.3}
\end{figure}



\subsection{从理论角度优化梯度压缩算法性能}

 除了用自适应策略强化算法鲁棒性，还有一些工作利用理论分析，提升算法在Non-IID场景下的收敛速率。他们通过对梯度压缩算法进行适当的改进，如添加偏移项，改变压缩规则等方法降低对数据异质性因子的依赖，从而在Non-IID场景下获得更快的收敛速率。有工作\cite{stich2020communication}比较了在两种适用于分布式机器学习的最常见的梯度压缩方法：Q-SGD和D-EF-SGD。该工作推导了这两种算法在Non-IID情况下的收敛速度，并得出结论：在Non-IID场景下，尽管D-EF-SGD仍然放大了数据异质性对收敛带来的负面影响，它的表现仍优于D-QSGD。还有工作\cite{hardthreshold}提出了硬阈值算法的稀疏化压缩策略（只发送绝对值高于恒定硬阈值的梯度），该工作构建了一个理论框架来推导不同稀疏化算法的收敛速度，他们基于这个框架推导出了硬阈值算法的的收敛速度，并证明了硬阈值压缩器对Non-IID场景更加稳健。
 

\section{数据量感知的算法}

目前仅有少数工作提出在分布式ML中进行数据量感知的节点选择\cite{amazon}。与传统的基于梯度的节点选择方法不同\cite{wu2022node}，这种选择方法纯粹基于节点数据量来选择参与该轮通信的节点，他们将重节点定义为那些数据集大小在前20$\%$的节点。他们的实验表明，在连续的FL中，如果给具有大数据量的节点以更高的选择概率，全局模型的质量就会提高。从数据大小选择节点是非常低成本和有效的。这项工作通过实验总结出结论，在Non-IID场景下，基于数据量的节点选择方法优于统一随机选择策略。他们的实验从另一个角度证明，不同尺寸的节点，他们的数据重要性是不同的，我们应该让大节点传递更多的有效信息，而不是采用同等对待的均匀压缩策略。

\iffalse
\section{基于梯度的节点选择}

如何在动态的网络环境和系统与数据的高异质性条件下合理地选择候选节点，是新型分布式机器学习的核心问题。合理选择候选节点可以有效提高收敛速度，提升模型质量。目前，针对数据异质性的节点选择主要是基于梯度或损失。服务器使用这些参数来确定客户是否偏离全局模型。如果一个客户端的迭代方向与全局模型相同，服务器会给它更高的权重或被选中的概率。这个领域最先进的工作是 \cite{wu2022node}。这项工作根据局部梯度和全局更新的内部乘积来调整选择节点的概率。他们提出了FedNPS，一个根据内积符号动态调整节点选择概率的框架。然而，基于梯度或损失的选择是非常昂贵的。有工作\cite{dupuy2022learnings}最近提出了一个简单的方法，根据数据大小选择节点。他们将重型设备定义为那些数据集大小在前20美元/%$的设备。他们的实验表明，在连续的FL中，如果给具有大数据量的节点以更高的选择概率，全局模型的质量就会提高。从数据大小选择节点是非常低成本和有效的。

\fi

上述（自适应）梯度压缩算法忽略了数据量的影响。也就是说， 各节点拥有的数据量的差异（反应在了各节点的训练权重差异）对其结果没有实际影响。如果对拥有大量数据集的节点使用与节点相同的激进压缩策略，训练中的关键信息就会丢失，从而导致严重的精度下降。我们提出的算法DAGC则充分考虑了节点尺寸差异的问题，它为有大量数据集的大节点设置了保守的压缩率，并为小节点分配了积极的压缩率。在固定和有限的通信条件下，DAGC实现了压缩比的最佳自适应调整。